\section{Positioning of this thesis}\label{section:background-positioning-of-thesis}

The preceding sections describe components that each have their own literature, but are usually treated separately: table, figure, and text extraction from PDFs, biomedical entity normalization, LLM-based extraction, natural-language inference, and cost-aware model routing. The position of this thesis is that an auditable clinical knowledge base cannot be built by any of these components alone, and should be integrated into a single pipeline, in which provenance is preserved, every artifact, and more importantly, every emitted rule can be traced back to its source, and finally, quality is evaluated together with cost over the whole pipeline. 

The system built in this thesis, therefore, consists of three parts. First, it uses a provenance-preserving document extraction stage, built on Docling, and extended with a hybrid table detector and post-processing stages. Second, it implements a grounded, atomic, and normalized rule-extraction pipeline, preserving source-faithfulness, and using scispaCy and UMLS with a frozen biomedical NLI backbone~\parencite{}. Third, it adds an agreement-based cascading pipeline~\parencite{} that is calibrated on a histopathology corpus. 

The main contribution of the thesis is the combination of all these components: the thesis measures what it costs to keep every extracted rule traceable to its evidence, where existing components are already sufficient off-the-shelf, and where additional stages in the pipeline, such as grounding, agreement, and escalation, are worth their cost. The results are reported and discussed in Chapters~\ref{chapter:Results} and~\ref{chapter:Discussion}.